% Part: methods
% Chapter: proofs
% Section: using-definitions

\documentclass[../../../include/open-logic-section]{subfiles}

\begin{document}

\olfileid{mth}{prf}{def}

\olsection{व्याख्यांचा वापर}

सिद्धतेत वापरल्या जाऊ शकणाऱ्या सर्व व्याख्या तुम्हाला परिचित असल्या पाहिजेत आणि त्यांचा योग्य वापर करता आला पाहिजे, असे आपण म्हटले होते. हा मुद्दा फार महत्त्वाचा आहे; त्याचा थोडा अधिक तपशील पाहू. गुणधर्म आणि संबंधांविषयी थोडक्यात बोलता यावे म्हणून व्याख्यांद्वारे त्यांना संक्षिप्त नावे दिली जातात. असे नाव म्हणजे \emph{व्याख्येय}; ते ज्याचा संक्षेप करते तो मजकूर म्हणजे \emph{व्याख्यापक}. सिद्धतेत अनेकदा व्याख्येय कसे मांडले होते याकडे परत जावे लागते. कारण सिद्धता पुढे नेण्यासाठी व्याख्यापकाची तार्किक रचना वापरावी लागते: व्याख्येय संज्ञा ज्याचा संक्षेप आहे तोच हा विस्तृत मजकूर असतो. व्याख्या उलगडल्यामुळे प्रत्यक्ष तार्किक कार्य जिथे होते तिथपर्यंत आपण पोहोचतो.

एका उदाहरणाने सुरुवात करू. पुढील विधान सिद्ध करायचे आहे असे समजा:

\begin{prop}
कोणत्याही $A$ आणि $B$ संचांसाठी, $A \cup B = B \cup A$.
\end{prop}

सिद्धतेची सुरुवात करण्यासाठीच दोन संच एकरूप असणे म्हणजे काय हे समजणे आवश्यक आहे. म्हणजे त्या समीकरणातील “$=$” चा संचांसाठी काय अर्थ आहे हे समजले पाहिजे. दोन संचांचे !!{element}s तेच असतील तेव्हा ते संच एकरूप असतात, अशी व्याख्या आहे. म्हणून उलगडायची व्याख्या ही आहे:

\begin{defn}
$A$ आणि $B$ संच \emph{एकरूप} आहेत, $A = B$, तेव्हाच आणि फक्त तेव्हा जेव्हा~$A$ मधील प्रत्येक !!{element} हा~$B$ मधील !!a{element} असतो आणि उलटही तसेच असते.
\end{defn}

या व्याख्येत $A$ आणि~$B$ ही स्वेच्छ संचांसाठी वापरलेली प्रतिनिधी नावे आहेत. ही व्याख्या ज्या गोष्टीची व्याख्या करते---म्हणजे \emph{व्याख्येय}---ती “$A = B$” ही अभिव्यक्ती आहे. $A = B$ कोणत्या अटीवर सत्य असते हे व्याख्या सांगते. ती अट---“~$A$ मधील प्रत्येक !!{element} हा~$B$ मधील !!a{element} असतो आणि उलटही तसेच असते”---म्हणजे \emph{व्याख्यापक}.\footnote{या विशिष्ट प्रकरणात---आणि येथे फार गोंधळ होऊ शकतो!---$A = B$ असते तेव्हा $A$ आणि $B$ हे प्रत्यक्षात एकच संच असतात, जरी त्यासाठी डावीकडे आणि उजवीकडे वेगवेगळी अक्षरे वापरली असली तरी. पण त्या संचाची ओळख करून देण्याचे मार्ग वेगळे असू शकतात; त्यामुळे व्याख्या निरर्थक पुनरुक्ती ठरत नाही.} ही अट पूर्ण होते तेव्हाच आणि फक्त तेव्हा (इंग्रजी संक्षेप “iff”) $A = B$ सत्य असते, असे ही व्याख्या सांगते.

व्याख्या वापरताना तिच्यातील $A$ आणि $B$ यांचा संबंध प्रत्यक्ष हाताळत असलेल्या प्रकरणाशी योग्य रीतीने जुळवला पाहिजे. आपल्या प्रकरणात $A \cup B = B \cup A$ सत्य असण्यासाठी प्रत्येक $z \in A \cup B$ हा $B \cup A$ मध्येही असला पाहिजे आणि उलटही तसेच असले पाहिजे. विधानातील $A \cup B$ ही अभिव्यक्ती व्याख्येतील~$A$ ची भूमिका बजावते आणि $B \cup A$ ही व्याख्येतील~$B$ ची. व्याख्येत आणि आपण सिद्ध करत असलेल्या विधानात $A$ आणि $B$ हीच नावे वेगवेगळ्या भूमिकांत येतात; म्हणून या दोन वापरांचा गोंधळ होणार नाही याची काळजी घ्यावी. उदाहरणार्थ,~$A$ मधील प्रत्येक !!{element} हा~$B$ मधील !!a{element} आहे आणि उलटही तसेच आहे, हे दाखवून विधान सिद्ध करता येईल असे समजणे चुकीचे होईल. त्यामुळे $A = B$ सिद्ध होईल; $A \cup B = B \cup A$ सिद्ध होणार नाही. शिवाय, $A$ आणि $B$ हे कोणतेही दोन संच असू शकतात. $A$ आणि~$B$ यांच्याविषयी काहीच गृहीत धरलेले नसेल, तर ते वेगळे संचही असू शकतात; म्हणून त्या चुकीच्या मार्गाने फार पुढे जाता येणार नाही.

सिद्धतेत संयोगासारख्या संचसैद्धान्तिक संकल्पना येतात. त्यामुळे सिद्धता कशी पुढे न्यायची हे समजण्यासाठी $\cup$ या चिन्हाचाही अर्थ समजला पाहिजे. कधी एखादी व्याख्या उलगडल्यावर आणखी व्याख्या उलगडाव्या लागतात. उदाहरणार्थ, $A \cup B$ ची व्याख्या $\Setabs{z}{z \in A \text{ किंवा } z \in B}$ अशी आहे. म्हणून $x \in A \cup B$ सिद्ध करायचे असेल, तर $\cup$ ची व्याख्या उलगडल्यावर $x \in \Setabs{z}{z \in A \text{ किंवा } z \in B}$ सिद्ध करावे लागेल असे दिसते. आता $x \in \Setabs{z}{\dots z\dots}$ हे $\dots x\dots$ असते तेव्हाच आणि फक्त तेव्हाच खरे असते, हेही आठवले पाहिजे. म्हणून $\Setabs{z}{\dots z \dots}$ या लेखनाची व्याख्या आणखी उलगडल्यावर असे दिसते की $x \in A$ किंवा $x \in B$ हे दाखवायचे आहे. त्यामुळे “$A \cup B$ मधील प्रत्येक !!{element} हा $B \cup A$ मधील !!a{element} देखील असतो” याचा प्रत्यक्ष अर्थ असा आहे: “प्रत्येक $x$ साठी, $x \in A$ किंवा $x \in B$ असेल, तर $x \in B$ किंवा $x \in A$ असते.” विधानातील व्याख्या पूर्ण उलगडल्यावर आपण हे दाखवायचे आहे असे दिसते:

\begin{prop}
कोणत्याही $A$ आणि $B$ संचांसाठी: (a) प्रत्येक $x$ साठी, $x \in A$ किंवा $x \in B$ असेल, तर $x \in B$ किंवा $x \in A$ असते; आणि (b) प्रत्येक $x$ साठी, $x \in B$ किंवा $x \in A$ असेल, तर $x \in A$ किंवा $x \in B$ असते.
\end{prop}

व्याख्या उलगडणे हा सिद्धता रचण्याचा आवश्यक भाग आहे, हा महत्त्वाचा मुद्दा आहे. हे योग्य रीतीने करणे कधी अवघड असते: व्याख्येतील चल आणि सिद्ध करावयाच्या दाव्यातील पदे यांतील फरक ओळखून त्यांची योग्य जुळवणी केली पाहिजे. यशस्वी होण्यासाठी प्रश्न नेमके काय विचारतो आणि त्यातील सर्व संज्ञांचा अर्थ काय आहे हे समजले पाहिजे. अनेकदा एकाहून अधिक व्याख्या उलगडाव्या लागतील. पुढे येणाऱ्या साध्या सिद्धतांत उत्तर जवळजवळ थेट व्याख्यांवरूनच मिळते. अर्थात, प्रत्येक वेळी हे इतके सोपे असेलच असे नाही.

\begin{prob}
$A \cap B \neq \emptyset$ सिद्ध करायला सांगितले आहे असे समजा. येथे येणाऱ्या सर्व व्याख्या उलगडा; म्हणजे “$\cap$”, “=” किंवा “$\emptyset$” यांचा उल्लेख न करता हे विधान पुन्हा मांडा.
\end{prob}

\end{document}
